\begin{pastexam}[label={px:ncku-108-prog-02}]{成大資訊聯招 108 學年度 第 2 題}
（計算）

Consider the following AOE network: (18\%)
\begin{center}
\begin{tikzpicture}[x=0.0145cm,y=-0.0145cm,>={Latex[length=2.2mm,width=1.6mm]},
  v/.style={circle,draw,fill=black!18,inner sep=0pt,minimum size=7mm,font=\itshape},
  l/.style={inner sep=1.5pt}]
\node[v] (v0) at (133,172) {$v_0$};
\node[v] (v1) at (302,70) {$v_1$};
\node[v] (v2) at (302,230) {$v_2$};
\node[v] (v3) at (213,340) {$v_3$};
\node[v] (v4) at (455,160) {$v_4$};
\node[v] (v5) at (455,340) {$v_5$};
\node[v] (v6) at (640,70) {$v_6$};
\node[v] (v7) at (587,215) {$v_7$};
\node[v] (v8) at (750,240) {$v_8$};
\node[left=2pt of v0] {start};
\node[right=5pt of v8] {finish};
\draw[->] (v0) -- (v1);
\draw[->] (v0) -- (v2);
\draw[->] (v0) -- (v3);
\draw[->] (v1) -- (v4);
\draw[->] (v2) -- (v4);
\draw[->] (v3) -- (v5);
\draw[->] (v4) -- (v6);
\draw[->] (v4) -- (v7);
\draw[->] (v5) -- (v7);
\draw[->] (v6) -- (v8);
\draw[->] (v7) -- (v8);
\draw[->] (v5) -- (v8);
\node[l] at (203,97) {$a_1=9$};
\node[l] at (236,170) {$a_2=4$};
\node[l] at (128,272) {$a_3=7$};
\node[l] at (405,75) {$a_4=2$};
\node[l] at (362,176) {$a_5=2$};
\node[l] at (333,320) {$a_6=4$};
\node[l] at (527,78) {$a_7=5$};
\node[l] at (555,145) {$a_8=7$};
\node[l] at (474,264) {$a_9=7$};
\node[l] at (753,112) {$a_{10}=1$};
\node[l] at (660,181) {$a_{11}=1$};
\node[l] at (650,325) {$a_{12}=3$};
\end{tikzpicture}
\end{center}
\begin{enumerate}[label=(\arabic*)]
\item Obtain $e(i)$ and $l(i)$ for all activity $i$. (7\%)
\item List all critical activities. (7\%)
\item List all critical paths. (4\%)
\end{enumerate}
\end{pastexam}

\begin{pastsoln}
記號照資結書第~22.5~節：事件的最早、最晚時間 $\mathit{ee}$、$\mathit{le}$，活動 $a_i = \langle j, k
\rangle$ 的 $e(i) = \mathit{ee}[j]$、$l(i) = \mathit{le}[k] - \mathit{dur}(a_i)$。 每條邊都從編號小的事件指向編
號大的，$v_0, v_1, \ldots, v_8$ 就是一個拓樸順序。

\textbf{往前算 $\mathit{ee}$、往回算 $\mathit{le}$：}

\begin{center}
\small
\begin{tabular}{cll}
\toprule
事件 & $\mathit{ee}$ & $\mathit{le}$ \\
\midrule
$v_0$ & $0$ & $\min(9 - 9,\ 9 - 4,\ 7 - 7) = 0$ \\
$v_1$ & $0 + 9 = 9$ & $11 - 2 = 9$ \\
$v_2$ & $0 + 4 = 4$ & $11 - 2 = 9$ \\
$v_3$ & $0 + 7 = 7$ & $11 - 4 = 7$ \\
$v_4$ & $\max(9 + 2,\ 4 + 2) = 11$ & $\min(18 - 5,\ 18 - 7) = 11$ \\
$v_5$ & $7 + 4 = 11$ & $\min(18 - 7,\ 19 - 3) = 11$ \\
$v_6$ & $11 + 5 = 16$ & $19 - 1 = 18$ \\
$v_7$ & $\max(11 + 7,\ 11 + 7) = 18$ & $19 - 1 = 18$ \\
$v_8$ & $\max(16 + 1,\ 18 + 1,\ 11 + 3) = 19$ & $19$（結束事件，$= \mathit{ee}[v_8]$） \\
\bottomrule
\end{tabular}
\end{center}

最短工期 $T = 19$。 開始事件的 $\mathit{le}$ 算回 $0$，驗算通過。

\textbf{(1) 每項活動的 $e$、$l$：}

\begin{center}
\small
\begin{tabular}{cccccc}
\toprule
活動 & 邊 & 工期 & $e$ & $l$ & $l - e$ \\
\midrule
$a_1$ & $\langle v_0, v_1 \rangle$ & 9 & 0 & $9 - 9 = 0$ & 0 \\
$a_2$ & $\langle v_0, v_2 \rangle$ & 4 & 0 & $9 - 4 = 5$ & 5 \\
$a_3$ & $\langle v_0, v_3 \rangle$ & 7 & 0 & $7 - 7 = 0$ & 0 \\
$a_4$ & $\langle v_1, v_4 \rangle$ & 2 & 9 & $11 - 2 = 9$ & 0 \\
$a_5$ & $\langle v_2, v_4 \rangle$ & 2 & 4 & $11 - 2 = 9$ & 5 \\
$a_6$ & $\langle v_3, v_5 \rangle$ & 4 & 7 & $11 - 4 = 7$ & 0 \\
$a_7$ & $\langle v_4, v_6 \rangle$ & 5 & 11 & $18 - 5 = 13$ & 2 \\
$a_8$ & $\langle v_4, v_7 \rangle$ & 7 & 11 & $18 - 7 = 11$ & 0 \\
$a_9$ & $\langle v_5, v_7 \rangle$ & 7 & 11 & $18 - 7 = 11$ & 0 \\
$a_{10}$ & $\langle v_6, v_8 \rangle$ & 1 & 16 & $19 - 1 = 18$ & 2 \\
$a_{11}$ & $\langle v_7, v_8 \rangle$ & 1 & 18 & $19 - 1 = 18$ & 0 \\
$a_{12}$ & $\langle v_5, v_8 \rangle$ & 3 & 11 & $19 - 3 = 16$ & 5 \\
\bottomrule
\end{tabular}
\end{center}

\textbf{(2)} 關鍵活動是 $l = e$ 的：$a_1, a_3, a_4, a_6, a_8, a_9, a_{11}$。

\textbf{(3)} 從 $v_0$ 出發只沿著關鍵活動走到 $v_8$（資結書定理~22.6：關鍵路徑就是長度 $T$ 的路
徑），有兩條：
\begin{align*}
&v_0 \xrightarrow{a_1} v_1 \xrightarrow{a_4} v_4 \xrightarrow{a_8} v_7 \xrightarrow{a_{11}} v_8
&& 9 + 2 + 7 + 1 = 19, \\
&v_0 \xrightarrow{a_3} v_3 \xrightarrow{a_6} v_5 \xrightarrow{a_9} v_7 \xrightarrow{a_{11}} v_8
&& 7 + 4 + 7 + 1 = 19.
\end{align*}
注意 $v_4$、$v_5$、$v_8$ 都是關鍵事件（$\mathit{ee} = \mathit{le}$），但 $a_{12} = \langle v_5, v_8 \rangle$ 有
$5$ 單位的寬裕，不在關鍵路徑上：兩端都是關鍵事件的活動不一定關鍵（資結書第~22.5~節）。
\end{pastsoln}
